\chapter{समानुपात और आँकड़े}\label{ch:g6:propdata}

तीन कॉपियाँ $6$ यूरो की हों, तो छह कॉपियाँ $12$ की: कॉपियाँ
दुगुनी, दाम दुगुना। ऐसा बरतने वाली राशियाँ \emph{\omterm{def:g6:propdata:def}{समानुपाती}} हैं
--- पूरे गणित के सबसे काम के विचारों में से एक, जिसे
\cref{ch:g7:prop} में और आगे बढ़ाया गया है। अध्याय के अंत में
सादे आँकड़ा-आलेख पढ़ना और बनाना है।

\section{समानुपाती राशियाँ}

\begin{definition}[समानुपात]\label{def:g6:propdata:def}
दो राशियाँ \emph{समानुपाती}\index{समानुपात} तब होती हैं जब एक के
मान दूसरी के मानों को हमेशा \emph{एक ही} संख्या से गुणा करके
मिलते हों; उस संख्या को \emph{समानुपात गुणांक} कहते
हैं।
\end{definition}

\begin{example}\label{ex:g6:propdata:table}
$2$-$2$ यूरो की कॉपियाँ:
\begin{center}
\renewcommand{\arraystretch}{1.3}
\begin{tabular}{l|cccc}
कॉपियाँ & $3$ & $5$ & $8$ & $12$ \\
\hline
दाम (यूरो) & $6$ & $10$ & $16$ & $24$ \\
\end{tabular}
\end{center}
हर दाम कॉपियों की संख्या $\times 2$ है: गुणांक $2$ है (एक का
दाम)। सारणी \omterm{def:g6:propdata:def}{समानुपाती} है या नहीं, इसकी झटपट जाँच: सारे
``स्तंभ-भागफल'' $\frac{6}{3}, \frac{10}{5}, \frac{16}{8},
\frac{24}{12}$ बराबर होने चाहिए --- यहाँ सब $2$ हैं।
\end{example}

\begin{example}[एक उलटा उदाहरण]\label{ex:g6:propdata:nonexample}
उम्र और ऊँचाई \omterm{def:g6:propdata:def}{समानुपाती} नहीं हैं: $12$ साल का बच्चा $6$ साल के
बच्चे से दुगुना लंबा नहीं होता। संख्याओं पर जाँचो: $12$ साल पर
$1.50$ m और $6$ साल पर $1.15$ m से \omterm{def:g3:division:remainder}{भागफल} $\frac{1.50}{12} = 0.125$
और $\frac{1.15}{6} \approx 0.19$ मिलते हैं --- बराबर नहीं।
\end{example}

\begin{method}[समानुपात की सारणी पूरी करना]\label{met:g6:propdata:complete}
तीन तरीक़े, जो चाहो चुनो:
\begin{enumerate}
  \item \emph{गुणांक}: किसी एक पूरे स्तंभ से गुणक निकालो, और उसे
        बाक़ी स्तंभों पर लगाओ;
  \item \emph{स्तंभ}: किसी स्तंभ को किसी संख्या से गुणा किया जा
        सकता है, या दो स्तंभ जोड़कर नया स्तंभ बनाया जा सकता है;
  \item \emph{एक पर लौटो}: पहले \emph{एक} वस्तु का मान निकालो,
        फिर गुणा करो।
\end{enumerate}
\end{method}

\begin{example}\label{ex:g6:propdata:methods}
पाँच एक जैसे मग $15$ यूरो के हैं; आठ मग कितने के होंगे?

\emph{एक पर लौटना}: एक मग $15 \div 5 = 3$ यूरो का है, इसलिए आठ
मग $8 \times 3 = 24$ यूरो के।

\emph{स्तंभ}: $8 = 5 + 3$; तीन मग $15$ का $\frac{3}{5}$, यानी
$9$ के हैं; इसलिए आठ मग $15 + 9 = 24$ यूरो के। वही उत्तर, जैसा
होना ही चाहिए।
\end{example}

\section{प्रतिशत}

\begin{definition}[प्रतिशत]\label{def:g6:propdata:percent}
``किसी राशि का $t\,\%$'' का मतलब है उसकी \omterm{def:g6:fractions:def}{भिन्न} $\frac{t}{100}$:
$60$ का $25\,\%$, $\frac{25}{100} \times 60 = 15$ है। प्रतिशत
गुणांक $\frac{t}{100}$ वाला \omterm{def:g6:propdata:def}{समानुपात} है।
\end{definition}

\begin{example}\label{ex:g6:propdata:percent}
$30$ विद्यार्थियों की कक्षा में $40\,\%$ लड़कियाँ हैं। लड़कियों की
संख्या, कदम-दर-कदम: $40\,\%$ यानी $\frac{40}{100} = 0.4$, और
$0.4 \times 30 = 12$ लड़कियाँ। काम के ठिकाने:
$50\,\%$ यानी \omterm{def:g3:division:half}{आधा}, $25\,\%$ यानी \omterm{def:g3:division:half}{चौथाई}, $10\,\%$ यानी दसवाँ भाग ---
तो $30$ का $10\,\%$, $3$ है, और $40\,\% = 4 \times 10\,\%$,
$4 \times 3 = 12$: वही उत्तर, मन-ही-मन निकला हुआ।
\end{example}

\section{आलेख पढ़ना और बनाना}

\begin{method}[सारणी या आलेख पढ़ना]\label{met:g6:propdata:read}
चित्र कुछ भी हो --- सारणी, पट्टी-आलेख, रेखा-आलेख:
\begin{enumerate}
  \item पहले शीर्षक और दोनों अक्षों के मात्रक पढ़ो;
  \item किसी सवाल का उत्तर देने के लिए सही पट्टी/स्तंभ ढूँढ़ो,
        फिर अंकन के सामने ध्यान से मान पढ़ो;
  \item तुलना के लिए पट्टियों की ऊँचाई देखो --- पर जाँचो कि अक्ष
        $0$ से शुरू होती है, नहीं तो चित्र धोखा दे सकता है।
\end{enumerate}
\end{method}

\begin{example}\label{ex:g6:propdata:chart}
नीचे का पट्टी-आलेख दिखाता है कि एक हफ़्ते में विद्यालय के
पुस्तकालय से कितनी किताबें ली गईं।

\begin{omfigure}
\begin{tikzpicture}
\begin{axis}[
  width=8.6cm, height=5.2cm,
  ybar, bar width=16pt,
  symbolic x coords={Mon,Tue,Wed,Thu,Fri},
  xtick=data, xticklabels={सोम,मंगल,बुध,गुरु,शुक्र},
  ymin=0, ymax=22,
  ytick={5,10,15,20},
  ylabel={ली गई किताबें},
  tick label style={font=\small},
  label style={font=\small}]
  \addplot[fill=omDef!30, draw=omDef] coordinates
    {(Mon,12) (Tue,8) (Wed,17) (Thu,6) (Fri,15)};
\end{axis}
\end{tikzpicture}

{\small एक पट्टी-आलेख: हर दिन के लिए एक पट्टी, ऊँचाई गिनती के बराबर।}
\end{omfigure}

पढ़ना: सबसे व्यस्त दिन बुधवार है ($17$ किताबें)। मंगलवार और
गुरुवार मिलकर $8 + 6 = 14$ किताबें बनते हैं --- अकेले शुक्रवार
($15$) से भी कम। हफ़्ते का कुल:
$12 + 8 + 17 + 6 + 15 = 58$ किताबें।
\end{example}

\begin{omfigure}
\begin{tikzpicture}
\begin{axis}[omaxis,
  width=8.2cm, height=5.4cm,
  xmin=0, xmax=9.4, ymin=0, ymax=28,
  xlabel={कॉपियाँ}, ylabel={दाम (यूरो)},
  xtick={1,2,3,4,5,6,7,8,9}, ytick={4,8,12,16,20,24}]
  \addplot[omDef, thick, domain=0:9] {3*x};
  \addplot[omDef, only marks, mark=*, mark size=1.5pt] coordinates
    {(1,3) (2,6) (3,9) (4,12) (5,15) (6,18) (7,21) (8,24)};
  \draw[dashed, omThm] (axis cs:6,0) -- (axis cs:6,18)
    -- (axis cs:0,18);
\end{axis}
\end{tikzpicture}

{\small \omterm{def:g6:propdata:def}{समानुपाती} राशियों का आलेख बहुत पहचाना-सा होता है: बिंदु
मूल बिंदु से होकर जाती एक सीधी रेखा पर आ जाते हैं (यहाँ
$3$-$3$ यूरो की कॉपियाँ; बिंदुदार रेखाएँ $6$ कॉपियों का दाम
पढ़ती हैं: $18$ यूरो)।}
\end{omfigure}

\section{अभ्यास}

\begin{exercise}[$\star$]\label{exo:g6:propdata:1}
क्या सारणी \omterm{def:g6:propdata:def}{समानुपाती} है? \omterm{def:g3:division:remainder}{भागफल} निकालकर पक्का करो।
\begin{center}
\renewcommand{\arraystretch}{1.2}
\begin{tabular}{l|ccc}
सेब (kg) & $2$ & $3$ & $5$ \\
\hline
दाम (यूरो) & $5$ & $7.5$ & $12.5$ \\
\end{tabular}
\qquad
\begin{tabular}{l|ccc}
उम्र (साल) & $2$ & $4$ & $8$ \\
\hline
ऊँचाई (cm) & $85$ & $103$ & $130$ \\
\end{tabular}
\end{center}
\end{exercise}

\begin{exercise}[$\star$]\label{exo:g6:propdata:2}
चार क्रोसाँ $4.80$ यूरो के हैं। एक क्रोसाँ का दाम निकालो,
फिर सात का।
\end{exercise}

\begin{exercise}[$\star$]\label{exo:g6:propdata:3}
\omterm{def:g6:propdata:def}{समानुपात} की सारणी पूरी करो (पहले गुणांक!):
\begin{center}
\renewcommand{\arraystretch}{1.2}
\begin{tabular}{l|cccc}
ईंधन (लीटर) & $10$ & $25$ & $?$ & $60$ \\
\hline
दाम (यूरो) & $18$ & $?$ & $72$ & $?$ \\
\end{tabular}
\end{center}
\end{exercise}

\begin{exercise}[$\star$]\label{exo:g6:propdata:4}
एक कार $100$ km पर $6$ L ईंधन लेती है। $250$ km के लिए कितना
ईंधन? $350$ km के लिए? $27$ L में वह कितनी दूर जा सकती है?
\end{exercise}

\begin{exercise}[$\star$]\label{exo:g6:propdata:5}
\cref{ex:g6:propdata:percent} के ठिकानों से मन-ही-मन निकालो:
$84$ का $50\,\%$; \quad $200$ का $25\,\%$; \quad $63$ का $10\,\%$;
\quad $70$ का $30\,\%$।
\end{exercise}

\begin{exercise}[$\star$]\label{exo:g6:propdata:6}
$450$ विद्यार्थियों के विद्यालय में $60\,\%$ भोजनालय में खाते
हैं। वे कितने विद्यार्थी हुए? कितने नहीं खाते?
\end{exercise}

\begin{exercise}[$\star$]\label{exo:g6:propdata:7}
\cref{ex:g6:propdata:chart} के पट्टी-आलेख से: किन दिनों में
$10$ से कम किताबें ली गईं? बुधवार को गुरुवार से कितनी ज़्यादा
किताबें ली गईं?
\end{exercise}

\begin{exercise}[$\star$]\label{exo:g6:propdata:8}
सोमवार से शुक्रवार तक दोपहर के तापमान $14$, $16$, $13$,
$17$, $20$ डिग्री थे। इन आँकड़ों का पट्टी-आलेख बनाओ (समझदारी से
अंकन चुनो), और सबसे गर्म दिन पढ़ो।
\end{exercise}

\begin{exercise}[$\star\star$]\label{exo:g6:propdata:9}
$6$ लोगों की एक पाकविधि में $450$ g आटा और $3$ अंडे लगते हैं।
उसे $10$ लोगों के लिए ढालो। (अंडों के लिए सोच लो, \omterm{def:g6:decimals:places}{दशमलव संख्या}
में अंडे लिखने से पहले!)
\end{exercise}

\begin{exercise}[$\star\star$]\label{exo:g6:propdata:10}
एक साइकिल सवार एक-सी चाल से $1$ घंटे में $24$ km चलती है।
\begin{enumerate}
  \item वह $2$ घंटे में कितनी दूर जाती है? आधे घंटे में?
        $1$ घंटा $30$ मिनट में?
  \item $60$ km के लिए उसे कितना समय चाहिए?
\end{enumerate}
\end{exercise}

\begin{exercise}[$\star\star$]\label{exo:g6:propdata:11}
एक दुकान ``हर चीज़ पर $20\,\%$ छूट'' देती है। इनके लिए छूट और नया
दाम निकालो: $15$ यूरो की गेंद; $40$ यूरो का रैकेट। क्या नया दाम
पुराने दाम का \omterm{def:g6:propdata:def}{समानुपाती} है? गुणांक क्या है?
\end{exercise}

\begin{exercise}[$\star\star\star$]\label{exo:g6:propdata:12}
एक ही ऊँचाई की दो मोमबत्तियाँ एक साथ जलाई जाती हैं। मोटी वाली
$6$ घंटे में पूरी जल जाती है, पतली वाली $4$ घंटे में, हर एक अपनी
एक-सी दर से। कितने समय बाद पतली मोमबत्ती ठीक मोटी वाली की आधी
ऊँचाई की रह जाती है? ($t$ घंटे बाद बची हुई ऊँचाइयाँ शुरू की
ऊँचाई की भिन्नों में लिखो।)
\end{exercise}

\section{समस्या: नक़्शे, नक़्शा-योजनाएँ, और आलेख से झूठ कैसे बोलें}

\begin{problem}[{सप्ताहांत समस्या --- नक़्शे की मापनी एक समानुपात
है: लंबाइयाँ गुणांक पर चलती हैं, क्षेत्रफल उसके वर्ग पर, और
ख़राब बने आलेख आँख को धोखा देते हैं}]
\label{pb:g6:propdata:1}
हर नक़्शा, फ़र्श-योजना और नमूना एक ही नियम मानता है: असली
लंबाइयाँ और बनी हुई लंबाइयाँ \emph{\omterm{def:g6:propdata:def}{समानुपाती}} होती हैं
(\cref{def:g6:propdata:def})। गुणांक का एक मशहूर नाम है ---
\emph{मापनी}\index{मापनी} --- और उसमें दो जाल छिपे हैं जो यह
समस्या जान-बूझकर बिछाती है: \omterm{def:g6:measure:area}{क्षेत्रफल} मापनी पर \emph{नहीं}
चलते, और प्रतिशत, आलेखों की तरह, पूरी तरह इस पर निर्भर करते हैं
कि उन्हें किसके सामने नापा गया है।

\medskip
\noindent\textbf{भाग I --- नक़्शा पढ़ना।}
पैदल यात्रा का एक नक़्शा $1:100\,000$ की मापनी का है: नक़्शे पर हर
सेंटीमीटर ज़मीन पर $100\,000$ सेंटीमीटर दिखाता है।
\begin{enumerate}
  \item $100\,000$ cm को किलोमीटर में बदलो। इस नक़्शे का $1$ cm
        km में क्या दिखाता है?
  \item नक़्शे पर दो गाँव $7.5$ cm दूर हैं; एक झील का \omterm{def:g5:solids:def}{किनारा}
        $12$ cm का वक्र है। असली दूरियाँ बताओ।
  \item एक बिलकुल सीधी पुरानी सड़क $20$ km लंबी है। वह नक़्शे पर
        कितनी लंबी है?
  \item दूसरा नक़्शा घोषित करता है: ``$1$ cm में $5$ km की दूरी''। उसकी
        मापनी $1: \,?$ के रूप में लिखो। उसी इलाक़े के लिए दोनों
        में से कौन-सा नक़्शा ज़्यादा बारीक़ी दिखाता है?
  \item पैदल यात्रा वाले नक़्शे के लिए एक छोटी सारणी बनाओ (नक़्शे
        की दूरी $1$, $2$, $7.5$ cm के सामने असली दूरी), और
        समझाओ कि मापनी ठीक \cref{def:g6:propdata:def} वाले अर्थ
        में \omterm{def:g6:propdata:def}{समानुपात} क्यों है। नक़्शे के सेंटीमीटरों को
        किलोमीटरों में बदलने वाला गुणांक क्या है?
\end{enumerate}

\medskip
\noindent\textbf{भाग II --- ज़ोई के कमरे की योजना।}
ज़ोई अपने कमरे --- $4$ m गुणा $3$ m के आयत --- की योजना
$1:50$ मापनी पर बनाती है।
\begin{enumerate}[resume]
  \item योजना पर कमरे के नाप क्या हैं?
  \item उसका पलंग $2$ m गुणा $0.9$ m का है, और दरवाज़े की चौड़ाई
        $80$ cm। तीनों नाप योजना पर बताओ।
  \item अब जाल। कमरे का असली \omterm{def:g6:measure:area}{क्षेत्रफल} m$^2$ में निकालो, फिर
        उसकी योजना का \omterm{def:g6:measure:area}{क्षेत्रफल} cm$^2$ में। असली \omterm{def:g6:measure:area}{क्षेत्रफल} को
        cm$^2$ में बदलो (\cref{def:g6:measure:area}) और योजना के
        \omterm{def:g6:measure:area}{क्षेत्रफल} से भाग दो: क्या \omterm{def:g3:division:remainder}{भागफल} $50$ है?
  \item जो संख्या मिली उसे समझाओ: $1:50$ की योजना पर काग़ज़ का हर
        वर्ग सेंटीमीटर $50$ cm गुणा $50$ cm के असली वर्ग को
        दिखाता है। वह कितने असली cm$^2$ हुए? नियम लिखो: लंबाइयाँ
        $50$ से भाग खाती हैं, तो \omterm{def:g6:measure:area}{क्षेत्रफल} किससे\,\dots
  \item एक गोल क़ालीन योजना पर $6$ cm$^2$ घेरता है। वह असल में
        कितना \omterm{def:g6:measure:area}{क्षेत्रफल} घेरता है, cm$^2$ में और फिर m$^2$ में?
\end{enumerate}

\medskip
\noindent\textbf{भाग III --- आलेख से (और प्रतिशत से) झूठ कैसे
बोलें।}
\begin{enumerate}[resume]
  \item एक दुकान ने शनिवार को $105$ यूरो का और रविवार को $110$
        यूरो का सामान बेचा। प्रबंधक ऐसा पट्टी-आलेख बनाता है
        जिसकी खड़ी अक्ष $100$ से शुरू होती है: शनिवार की पट्टी
        $5$ छोटी इकाई ऊपर जाती है, रविवार की $10$। चित्र रविवार
        के बारे में क्या सुझाता है, और सच क्या है? (रविवार की
        बढ़त शनिवार के प्रतिशत में निकालो, नज़दीकी प्रतिशत तक।)
        प्रबंधक ने \cref{met:g6:propdata:read} की कौन-सी
        चेतावनी अनसुनी की?
  \item ईमानदार आलेख दोबारा बनाओ (या बताओ), जिसकी अक्ष $0$ से
        शुरू हो: अब दोनों पट्टियाँ कैसी लगती हैं?
  \item किसी संग्रह में टिकट $40$ से बढ़कर $60$ हो जाते हैं:
        बढ़त को शुरुआती मान के प्रतिशत में निकालो। फिर वह $60$
        से घटकर $40$ पर आ जाता है: घटत को \emph{उसके} शुरुआती
        मान के प्रतिशत में निकालो। उन्हीं $20$ टिकटों के लिए
        दोनों उत्तर अलग क्यों हैं?
  \item छूट का मौसम: $100$ यूरो के कोट पर ``$20\,\%$ छूट'', और
        भुगतान पर घटे हुए दाम पर ``$10\,\%$ और छूट''। आख़िरी दाम
        कदम-दर-कदम निकालो। क्या कुल छूट $30\,\%$ है? ख़रीदार को
        समझाओ कि वह असल में कितनी है।
  \item आख़िरी सवाल, फिर प्रश्न 4 वाले नक़्शे पर, जहाँ $1$ cm
        में $5$ km हैं: एक जंगल उस नक़्शे पर $3$ cm$^2$ घेरता है। उसका
        असली \omterm{def:g6:measure:area}{क्षेत्रफल} km$^2$ में क्या है? इस पूरी समस्या का
        नियम लिखकर ख़त्म करो: लंबाइयाँ मापनी पर चलती हैं,
        \omterm{def:g6:measure:area}{क्षेत्रफल} उसके\,\dots
\end{enumerate}
\end{problem}
